(1) 弦点珠两侧的横向张力合力为 \(T[y_x(0^+,t)-y_x(0^-,t)]\)。点珠满足
\[
M y_{tt}(0,t)=T[y_x(0^+,t)-y_x(0^-,t)],
\]
同时位移连续 \(y(0^-,t)=y(0^+,t)\)。此跃变条件必须保留点质量，否则误以为斜率也连续会错误预测没有散射。

(2) 采用 \(e^{-i\omega t}\) 时间惯例，
\[
y_L=e^{ikx}+re^{-ikx},\qquad y_R=te^{ikx}
\]
表示各侧复振幅。位移连续得 \(t=1+r\)。导数跃变为 \(ik t-ik(1-r)=2ikr\)，代入点珠方程
\[
-M\omega^2t=2iTk r.
\]
令 \(Z=Tk/\omega=\sqrt{T\mu}\)，解得
\[
r=-\frac{M\omega}{M\omega+2iZ},\qquad
t=\frac{2iZ}{M\omega+2iZ}.
\]
它们还满足 \(t=1+r\)，可作代数核验。

(3) 两侧弦张力与波速相同，所以单色行波能流与振幅模平方成正比。故
\[
\mathcal R=|r|^2=\frac{M^2\omega^2}{M^2\omega^2+4Z^2},\quad
\mathcal T=|t|^2=\frac{4Z^2}{M^2\omega^2+4Z^2}.
\]
相加严格等于1，证明点珠无耗散，只将入射能量重新分成左行与右行波。虽然点珠暂时储存机械动能，周期平均后的能量通量仍平衡。

(4) 当 M→0，反射 \(\mathcal R\to0\)，恢复均匀弦；M→∞ 或固定 M 而 ω→∞，\(\mathcal R\to1\)，点珠几乎成为固定端。若 ω→0 且 M 有限，\(\mathcal T\to1\)，足够慢的波能让点珠跟随而少有散射。关键无量纲参数为 \(M\omega/(2Z)\)，因此这不是通过更换粒子质量数字产生的简单重复题。
